\section{Why minimizers, and why does this chapter exist?}
\label{sec:why-minimizers}

Chapter~\ref{r:{biology}} introduced the seed--chain--extend
pipeline used by every modern aligner.  In that pipeline, the
\emph{minimizer collector} is the only stage whose runtime cost is
strictly linear in the input length and not bounded by the alignment
depth: it must visit every base of every sequence it processes,
turning each $k$-mer into a hash, and selecting the
window-minimum of every $w$ consecutive $k$-mers as the
\emph{minimizer} for that window.

For our user-stated benchmarking goal -- pairs of ${\sim}10^6$ bp
sequences aligned against ${\sim}10^7$ bp references at PacBio HiFi
error rates -- the minimizer step accounts for roughly 8\,\% of the
total wall-clock per pair on the scalar pipeline, and up to 25\,\%
when a naive SIMD port has been written. That is the gap this
chapter is about.  Concretely:

\begin{itemize}
\item \textbf{Goal 1} (already met before this chapter starts).
  bioseq's \emph{scalar} pipeline must be within $\pm 10\,\%$ of
  minimap2's wall-clock on the long-sequence regime.  Met by the
  per-seed-pair gap-fill rewrite documented in V3 of the engineering
  report~\cite{bioseq_eng_report_v1}.
\item \textbf{Goal 2} (the topic of this chapter).
  bioseq's \emph{SIMD} pipeline must be at least 10\,\% faster than
  the scalar pipeline.  After a year of incremental work we sit
  marginally outside the goal on every length bucket: $-8.4\,\%$ at
  52\,kbp, $-8.8\,\%$ at 903\,kbp, $-4.9\,\%$ at 5\,Mbp (means of three
  runs).  The variance across runs occasionally clears the bar (V12
  individual runs at $-9.7\,\%$ to $-12\,\%$ have been seen on a
  quiescent host), but the mean reported here is the load-bearing
  number.  This chapter explains why every percentage point was so
  hard-won.
\end{itemize}

Figure~\ref{fig:headline-result} states the bottom line: across our
five sequence-length buckets, the V13 collector
(\texttt{collect\_minimizers\_two\_stack\_nthash\_v3}, the
production minimizer at the time of writing) is the fastest of the
SIMD configurations, and at every length except the 5\,Mbp XL bucket
it is also faster than the scalar baseline by 5--9\,\%.

\begin{figure}[ht]
  \centering
  \begin{CacheMeCode}{python, custom include command={\input{\robExtAddCachePathAndName{\robExtFinalHash.pgf}}}, set placeholder eval={__LINEWIDTH__}{\lenToCmNoUnit[in]{\linewidth}}}
    import matplotlib.pyplot as plt
    import matplotlib
    import numpy as np
    matplotlib.use("pgf")
    matplotlib.rcParams.update({
      "font.family": "serif",
      "font.serif": [],
      "text.usetex": True,
    })
    fig, ax = plt.subplots()
    fig.set_size_inches(__LINEWIDTH__, 0.55*__LINEWIDTH__, forward=True)
    buckets = ['47 kbp', '49 kbp', '52 kbp', '903 kbp', '5.0 Mbp']
    scalar = [14.56, 14.94, 15.09, 188.93, 1171.0]
    simd_v6 = [15.21, 15.90, 16.05, 201.68, 1207.99]
    v8_two_stack = [14.64, 15.07, 15.24, 183.77, 1131.87]
    v12_v13 = [13.78, 14.02, 13.82, 172.38, 1113.3]
    minimap2 = [20.11, 20.10, 20.13, 190.95, 1060.13]
    x = np.arange(len(buckets))
    width = 0.16
    ax.bar(x - 2*width, minimap2, width, label='minimap2',
           color='#000000')
    ax.bar(x - width,   scalar,    width, label='scalar (baseline)',
           color='#909090')
    ax.bar(x,           simd_v6,   width, label='V6 SIMD (rescan)',
           color='#1f77b4')
    ax.bar(x + width,   v8_two_stack, width, label='V8 two-stack',
           color='#2ca02c')
    ax.bar(x + 2*width, v12_v13,   width, label='V13 (production)',
           color='#d62728')
    ax.set_xticks(x)
    ax.set_xticklabels(buckets)
    ax.set_yscale('log')
    ax.set_ylabel('Wall-clock per pair (ms, log scale)', fontsize=8)
    ax.set_title('End-to-end pipeline wall-clock per length bucket', fontsize=9)
    ax.tick_params(axis='both', labelsize=7)
    ax.legend(fontsize=7, frameon=False, loc='upper left', ncol=2)
    ax.grid(True, alpha=0.3, axis='y', which='both')
    print(get_filename_from_extension(".pgf"))
    plt.savefig(get_filename_from_extension(".pgf"), bbox_inches="tight")
  \end{CacheMeCode}
  \caption{End-to-end pipeline wall-clock per pair, length bucket on
    the X axis (log Y). \textsc{Scalar} is the baseline, \textsc{V6}
    is the naive SIMDification of minimap2's rescan loop,
    \textsc{V8} is the two-stack port from the SIMD-Minimizers paper
    of Groot Koerkamp \&
    Martayan~\cite{groot_koerkamp_simd_minimizers_2025}, and
    \textsc{V13} is the production state after dual-stream
    interleaving, forced inlining, and the
    \texttt{mm\_est\_err} sort-free prefix-count rewrite.  The
    minimap2 reference is included for orientation. On every bucket
    except XL, V13 is the fastest of all five rows.}
  \label{fig:headline-result}
\end{figure}

\section{What a minimizer collector actually does}
\label{sec:what-is-minimizer}

Before we dive into the SIMD work, it pays to spend a page or two on
what the data looks like and how the original scalar algorithm gets
the answer.  This section is aimed at a reader who has seen the
phrase ``minimizer'' in passing -- e.g.\ from a Kraken or a minimap2
paper -- but has never written or read minimizer code.

Conceptually, a minimizer collector is a function
\[
\mathrm{collect\_minimizers}\colon \Sigma^* \times \mathbb{N} \times \mathbb{N}
\;\to\; \mathcal{P}\!\big(\;(\text{hash},\,\text{position},\,\text{strand})\;\big)
\]
that takes a sequence over the four-letter DNA alphabet, plus
parameters $w$ (window size) and $k$ ($k$-mer length), and returns a
set of triples that uniquely identify selected $k$-mers in the input.
The triples are what downstream stages -- chaining, error estimation,
seed-hit collection -- consume.

\begin{figure}[ht]
  \centering
  \begin{tikzpicture}[x=0.4cm, y=0.4cm, >=Latex]
    \foreach \i/\letter in {0/A,1/C,2/T,3/G,4/T,5/A,6/G,7/T,8/A,9/C,10/G,11/T,12/T,13/A}
        \node[minimum width=0.4cm,inner sep=0pt,font=\ttfamily\small] at (\i,0) {\letter};

    \foreach \i/\hashval in {0/13,1/20,2/27,3/34,4/41}
    {
      \pgfmathsetmacro{\xs}{\i - 0.2}
      \pgfmathsetmacro{\xe}{\i + 3.8}
      \pgfmathsetmacro{\yp}{-1.0 - \i*0.8}
      \draw[rounded corners=1pt, color=blue!50!gray, thick]
        (\xs, \yp) rectangle (\xe, \yp+0.6);
      \node[font=\tiny, color=blue!50!gray, anchor=west] at (\xe + 0.1, \yp+0.3)
        {kmer$_\i$ (h=\hashval)};
    }
    \node[font=\tiny, anchor=west, color=blue!50!gray]
      at (0, -5.6) {(kmers $k=4$ shown for first 5 positions)};

    \draw[rounded corners=1pt, color=ForestGreen, thick]
      (-0.3, -7.0) rectangle (8.3, -6.4);
    \node[font=\scriptsize, color=ForestGreen, anchor=west]
      at (-0.3, -7.6) {window of $w=4$ consecutive $k$-mers ($w+k-1=7$ bases wide)};
  \end{tikzpicture}
  \caption{The data the collector operates on. A sequence of 14
    bases produces overlapping $k=4$-mers
    (kmer$_0$, kmer$_1$, \ldots). The collector slides a
    window of $w=4$ consecutive $k$-mers (so $w+k-1=7$ bases wide)
    across the sequence and emits, for each window, the position of
    the $k$-mer with the smallest hash. A naive output would have
    one entry per window; a smart implementation deduplicates
    adjacent windows that share the same minimizer position.}
  \label{fig:minimizer-data}
\end{figure}

The intuition for \emph{why} a minimizer is a useful seed is in the
seed-stage discussion of Chapter~\ref{r:{biology}}:
the same $k$-mer is selected as the minimizer of every
overlapping window that contains it as the smallest hash, so two
sequences that share a contiguous identical region are very likely
to share at least one minimizer position there. The seed
collection then becomes a hash-table intersection rather than a
substring search, which is the entire reason any of this is fast.

\subsection{The minimap2 hash}
\label{sec:hash64}

Random minimizers depend critically on the hash function used to
order $k$-mers.  Three properties matter:

\begin{enumerate}
\item \textbf{Avalanche.} A single bit flip in the $k$-mer should
  change roughly half the bits of the hash output -- otherwise
  GC-rich and AT-rich windows occupy different
  hash-value subspaces and the minimizer set becomes biased.
\item \textbf{Reversibility.} The hash should be a
  bijection on its $\mathrm{2k}$-bit input domain, so that no two
  $k$-mers collide.  This is implementable with a sequence of
  reversible mixing operations
  (XOR, addition, multiplication by an odd constant, shifts).
\item \textbf{Cheap to evaluate.} The minimizer collector evaluates
  the hash once per $k$-mer position, so on a 1\,Mbp pair it runs
  the hash $2 \cdot 10^6$ times. Each cycle counts.
\end{enumerate}

minimap2 uses a hand-tuned 64-bit integer mix attributed to Thomas
Wang~\cite{wang_inthash}, applied to the bit-packed
$2k$-bit forward $k$-mer encoding (or to the reverse-complement
$k$-mer, whichever is lexicographically smaller -- this is the
\emph{canonical-strand} optimisation):

\begin{verbatim}
static inline uint64_t hash64(uint64_t key, uint64_t mask) {
    key = (~key + (key << 21)) & mask;        // key = key * 2097151 - 1
    key = key ^ (key >> 24);
    key = ((key + (key << 3)) + (key << 8)) & mask;   // * 265
    key = key ^ (key >> 14);
    key = ((key + (key << 2)) + (key << 4)) & mask;   // * 21
    key = key ^ (key >> 28);
    key = (key + (key << 31)) & mask;
    return key;
}
\end{verbatim}

Each pass alternates a multiplication-by-odd-constant (implemented as
a few \texttt{shift\,+\,add} pairs to be portable to
HP\,9000-class hardware where multiplication was slow in 1997) with
an XOR with a shifted version of itself.  The composite is
reversible (each stage is reversible) and has good
avalanche~\cite{wang_inthash}, but the dependency chain is
\textbf{seven stages deep}.  When the loop is in scalar code, the
out-of-order engine schedules these against the rest of the body and
hides the latency.  In the SIMD path the seven stages become a
seven-deep fused dependency chain on a single SIMD register, and they
end up dominating the inner loop -- a fact that drives most of the
optimisation work later in the chapter
(see~\S\ref{sec:v8-defer-scramble} for the eventual fix).

\subsection{The original minimap2 algorithm}

\texttt{mm\_sketch} in
\path{bench/tool_minimap2/minimap2/sketch.c} is the canonical scalar
minimizer collector.  Algorithm~\ref{alg:mm-sketch} captures the hot
path; for clarity we have stripped the homopolymer compression
branch and the bookkeeping for tied hashes.

\begin{algorithm}[H]
\caption{\texttt{mm\_sketch} (simplified; minimap2 \texttt{sketch.c})}
\label{alg:mm-sketch}
\small
\begin{algorithmic}[1]
\Require sequence $s$, window size $w$, $k$-mer length $k$
\State $\textit{buf}[0..w-1] \gets \infty$ \Comment{circular buffer of recent $k$-mer info}
\State $\textit{min} \gets \infty$, $\textit{min\_pos} \gets -1$
\State $\textit{kmer}[0] \gets 0$, $\textit{kmer}[1] \gets 0$ \Comment{forward / RC}
\State $\textit{shift1} \gets 2(k-1)$
\For{$i \gets 0$ to $|s|-1$}
    \State $c \gets \textit{seq\_nt4\_table}[s[i]]$ \Comment{ACGT $\to$ 0..3}
    \State $\textit{kmer}[0] \gets ((\textit{kmer}[0] \ll 2)\;|\;c)\,\&\,\textit{mask}$ \Comment{forward}
    \State $\textit{kmer}[1] \gets (\textit{kmer}[1] \gg 2)\;|\;((3 \oplus c) \ll \textit{shift1})$ \Comment{RC}
    \If{$\textit{kmer}[0] = \textit{kmer}[1]$} \Comment{palindrome}
        \State \textbf{continue} \Comment{strand undefined; skip}
    \EndIf
    \State $z \gets \textit{kmer}[0] < \textit{kmer}[1] \;?\; 0 : 1$ \Comment{strand}
    \State $\textit{info}.x \gets \textit{hash64}(\textit{kmer}[z]) \ll 8 \;|\; \textit{kmer\_span}$
    \State $\textit{info}.y \gets (i \ll 1)\;|\;z$
    \State $\textit{buf}[\textit{buf\_pos}] \gets \textit{info}$
    \If{$\textit{info}.x \le \textit{min}.x$}                                  \Comment{new minimum}
      \State \textbf{emit} $\textit{min}$ if armed
      \State $\textit{min} \gets \textit{info}$, $\textit{min\_pos} \gets \textit{buf\_pos}$
    \ElsIf{$\textit{buf\_pos} = \textit{min\_pos}$}                          \Comment{old min slid out}
      \State \textbf{emit} $\textit{min}$
      \State $\textit{min} \gets \infty$
      \For{$j \gets \textit{buf\_pos}+1$ \textbf{to} $w-1$ \textbf{then} $0$ \textbf{to} $\textit{buf\_pos}$}        \Comment{rescan window}
        \If{$\textit{min}.x \ge \textit{buf}[j].x$}
          \State $\textit{min} \gets \textit{buf}[j]$, $\textit{min\_pos} \gets j$
        \EndIf
      \EndFor
    \EndIf
    \State $\textit{buf\_pos} \gets (\textit{buf\_pos} + 1) \bmod w$
\EndFor
\State \textbf{emit} $\textit{min}$ at end of stream
\end{algorithmic}
\end{algorithm}

The design of Algorithm~\ref{alg:mm-sketch} is the
\emph{rescan} method described in Section~3.3 of the SIMD-Minimizers
paper~\cite{groot_koerkamp_simd_minimizers_2025}.
The salient property is that the loop has \emph{exactly one}
data-dependent branch on the hot path: the rescan branch in line~17
fires only when the current minimum slides out of the window.  In
expectation, the branch fires every $w$ iterations, so for the
default $w=10$ minimap2 has a $\sim 10\,\%$ branch rate that
modern out-of-order hardware predicts well.  Combined with the fact
that scalar code can keep all of $\textit{min}, \textit{kmer}[0],
\textit{kmer}[1], \textit{buf\_pos}, \textit{min\_pos}$ in
registers and let the OoO engine hide the seven-stage hash latency,
this loop achieves $\sim 12$\,ms per 1\,Mbp pair on an Apple\,M1
performance core.  That is the bar.

\section{Round 1 -- the naive SIMDification (V6)}
\label{sec:v6-rescan}

Our first SIMD port -- \texttt{collect\_minimizers\_simd} in
\path{bioseq/backend/sequences/kmers/index.cppm} -- was a
literal SIMDification of Algorithm~\ref{alg:mm-sketch}.  We split
the input sequence into $L$ equal-length chunks and ran the rescan
loop in parallel on $L$ lanes of a SIMD register.  On 128-bit Apple
NEON with \texttt{uint32\_t} storage, $L = 4$.

\begin{figure}[ht]
  \centering
  \begin{tikzpicture}[scale=0.55, transform shape, node distance=0.5cm, >=Latex]
    \node[align=center, font=\small] at (0,1.5)
      {Input sequence (1\,Mbp), split into $L=4$ equal chunks of $\sim 250$\,kbp};
    \draw[thick, fill=blue!10] (-5,0) rectangle (-2.5,1) node[midway, font=\small\ttfamily] {chunk 0};
    \draw[thick, fill=green!10] (-2.5,0) rectangle (0,1) node[midway, font=\small\ttfamily] {chunk 1};
    \draw[thick, fill=orange!10] (0,0) rectangle (2.5,1) node[midway, font=\small\ttfamily] {chunk 2};
    \draw[thick, fill=red!10] (2.5,0) rectangle (5,1) node[midway, font=\small\ttfamily] {chunk 3};

    \draw[->, thick] (-3.75, 0) -- (-3.75, -0.8);
    \draw[->, thick] (-1.25, 0) -- (-1.25, -0.8);
    \draw[->, thick] (1.25,  0) -- (1.25,  -0.8);
    \draw[->, thick] (3.75,  0) -- (3.75,  -0.8);

    \node[align=center, font=\small] at (0, -1.4)
      {One SIMD register holds one $k$-mer hash from each chunk:};

    \draw[thick] (-2.5,-2.5) rectangle (2.5,-1.7);
    \draw[thick] (-1.25,-2.5) -- (-1.25,-1.7);
    \draw[thick] (0,   -2.5) -- (0,   -1.7);
    \draw[thick] (1.25,-2.5) -- (1.25,-1.7);
    \node[font=\small\ttfamily] at (-1.875, -2.1) {h$_0$};
    \node[font=\small\ttfamily] at (-0.625, -2.1) {h$_1$};
    \node[font=\small\ttfamily] at (0.625,  -2.1) {h$_2$};
    \node[font=\small\ttfamily] at (1.875,  -2.1) {h$_3$};

    \node[font=\scriptsize] at (-3.5, -2.1) {SIMD reg:};
  \end{tikzpicture}
  \caption{Lane-parallel SIMD layout. The sequence is split into
    $L=4$ chunks; the SIMD register holds one $k$-mer hash from each
    chunk in parallel. Adjacent chunks overlap by $w + k - 1$ bases
    so that every window is fully contained in at least one chunk.}
  \label{fig:simd-chunks}
\end{figure}

The naive port retained the rescan structure of
Algorithm~\ref{alg:mm-sketch} but replaced every scalar comparison
with a SIMD comparison.  The rescan branch -- which fired
$\sim 10\,\%$ of the time in scalar code -- now fires whenever
\emph{any} of the four lanes sees its saved minimum slide out of
the window.  The probability that at least one of four lanes
triggers a rescan is $1 - (1 - 0.1)^4 \approx 34\,\%$, so we now
take a data-dependent branch one iteration in three.  Worse, the
rescan body itself becomes a $w$-iteration scalar loop nested
inside the SIMD body, because the rescan has to figure out
\emph{which} lane is being rescanned and keep the others
unchanged.  Modern branch predictors do not enjoy this combination.

The instrumented per-pair phase breakdown made the regression
unambiguous (Table~\ref{tab:v6-phase}): the scalar
\texttt{mm\_sketch} took 12.5\,ms on a 903\,kbp pair, while the
``vectorised'' \texttt{mm\_sketch\_simd} took 24\,ms -- a 2$\times$
slowdown despite running on the same hardware with four lanes of
parallelism.

\begin{table}[ht]
  \centering
  \begin{tabular}{lrrr}
    \toprule
    Phase & scalar (ms) & V6 SIMD (ms) & V6 over scalar \\
    \midrule
    \texttt{mm\_sketch}      & 12.50 & 14.94 & $+19\,\%$ \\
    \texttt{mm\_sketch\_simd} & 0    & 23.88 & --- \\
    \texttt{target\_index\_build} & 16.0 & 14.3 & $-11\,\%$ \\
    \texttt{collect\_seed\_hits} & 6.5  & 11.0 & $+69\,\%$ \\
    \texttt{mg\_lchain\_dp}   & 16.0 & 0.0  & --- (replaced) \\
    \texttt{mg\_lchain\_simd} & 0.0  & 8.9  & --- \\
    \texttt{mm\_est\_err}     & 0.7  & 7.2  & $+929\,\%$ \\
    \texttt{mm\_align\_skel.}  & 87.5 & 86.4 & $-1\,\%$ \\
    \midrule
    Pipeline total per pair & 158.4 & 170.7 & $+8\,\%$ \\
    \bottomrule
  \end{tabular}
  \caption{Phase breakdown at 903\,kbp, V3 baseline -- pre-SIMD-Minimizers
    work. The naive SIMD path (V6) is $+8\,\%$ \emph{slower} than
    scalar despite four lanes of parallelism. The slowdown is split
    between three regressions: (a) the SIMD minimizer's
    \texttt{mm\_sketch\_simd} is itself slower than scalar
    \texttt{mm\_sketch}; (b) the SIMD chainer feeds anchors into
    \texttt{collect\_seed\_hits} that hit a worse code path; and
    (c) the SIMD chain anchors are not cleanly indexable into the
    sorted query-minimizer array, so \texttt{mm\_est\_err} explodes.
    The V13 work mostly fixes (a) and (c).}
  \label{tab:v6-phase}
\end{table}

The post-mortem on V6 reads as a list of mistakes, each of which has
a corresponding lesson encoded into the next round:

\begin{enumerate}
\item \textbf{The rescan branch is more frequent in SIMD than in
  scalar.}  Independent lanes mean independent ``minimum slid out''
  events, and the rate compounds.  Fix: move to a data-independent
  algorithm (the two-stacks method, \S\ref{sec:v8-two-stack}).

\item \textbf{The seven-stage hash mix dominates the inner loop.}
  In scalar code the OoO engine hides the seven-cycle dependency
  chain behind the rest of the loop body.  In SIMD code there is no
  ``rest of the loop'' to hide it behind: the chain is the loop.
  Fix: defer the scramble to emit time, when only $\sim 1/w$ of
  positions need it (\S\ref{sec:v8-defer-scramble}).

\item \textbf{The \texttt{submit\_results} lambda fires per emit
  position}, and emits via a $w$-iteration per-lane loop because
  emissions on different lanes correspond to different output
  positions.  Each invocation costs roughly the same as a
  per-iteration call.  At $\sim 1/w$ emit rate this contributes
  $\sim 5$\,ms per Mbp.  Fix: hand-inline the per-step body and
  promote the emission loop into the inner frame (V9, V12).

\item \textbf{The hasher is parameterised on a runtime
  \texttt{rev\_shift = 2(k-1)}.}  When $k$ is a runtime parameter,
  the SIMD shift instruction has to load the shift count from a
  register every iteration.  Constant-amount NEON shifts (\texttt{shl.4s
  vN, vM, \#28}) are single-cycle, runtime shifts are slower.  Fix:
  add a \texttt{StaticK = 15} template parameter so that the
  minimap2-default $k$ produces a \texttt{constexpr}
  \texttt{rev\_shift} (V8 / V12).
\end{enumerate}

The downstream regressions on \texttt{collect\_seed\_hits} and
\texttt{mm\_est\_err} are not fundamentally about the minimizer
collector -- they are about the SIMD chainer producing a slightly
different anchor distribution than scalar -- but they are
\emph{visible} in the V6 phase breakdown because the SIMD chainer is
co-enabled with the SIMD minimizer in the bench.  The fixes for
those (the \texttt{mm\_est\_err} sort-free rewrite and the
divergence-fix for reverse-strand chains) eventually land in
V13~(\S\ref{sec:v13-est-err}) and reduce \texttt{mm\_est\_err} from
7\,ms to 0.85\,ms.

\section{Round 2 -- two-stacks port (V8)}
\label{sec:v8-two-stack}

The single most important external input to this work was the
SIMD-Minimizers paper of Groot Koerkamp \&
Martayan~\cite{groot_koerkamp_simd_minimizers_2025}.  Its central
insight is that the rescan method is fundamentally
\emph{data-dependent}, and that any data-dependent branch costs more
in SIMD than in scalar because lanes' branch targets are tied to
each other.  The paper proposes replacing the rescan with the
\emph{two-stacks} sliding-window-minimum algorithm of Hirzel,
Schneider \& Tangwongsan~\cite{hirzel_two_stacks_2017}, which has
\emph{zero} data-dependent branches -- only the once-per-$w$-iters
``rebuild the suffix-min array'' uniform branch.

\subsection{The two-stacks idea}

The two-stacks algorithm splits the stream of $k$-mer hashes into
fixed-size blocks of $w$ consecutive hashes each, computes the
prefix-min and suffix-min of each block in $O(w)$, and uses the fact
that any window of $w$ hashes spans (at most) one suffix of a
finished block plus a prefix of the current block.  The window
minimum is then the minimum of the corresponding suffix and prefix
minima -- both available in $O(1)$ from the precomputed arrays.

\begin{figure}[ht]
  \centering
  \begin{tikzpicture}[x=0.42cm, y=0.42cm, >=Latex]
    \foreach \i/\val in {0/9,1/4,2/7,3/2,4/8,5/5,6/3,7/6,8/1,9/9,10/4,11/7,12/2,13/8,14/5,15/3}
        \node[draw, minimum width=0.4cm, inner sep=0pt, font=\scriptsize\ttfamily] at (\i,3) {\val};

    \foreach \i in {0,4,8,12}
    {
      \pgfmathsetmacro{\xs}{\i - 0.25}
      \pgfmathsetmacro{\xe}{\i + 3.25}
      \draw[rounded corners=1pt, color=blue!50!black, ultra thick]
        (\xs, 2.5) rectangle (\xe, 3.5);
    }
    \node[anchor=east, font=\scriptsize, color=blue!50!black] at (-0.5,3) {hashes:};

    \foreach \i/\val in {0/9,1/4,2/4,3/2,4/8,5/5,6/3,7/3,8/1,9/1,10/1,11/1,12/2,13/2,14/2,15/2}
        \node[draw, minimum width=0.4cm, inner sep=0pt, font=\scriptsize\ttfamily, fill=green!10] at (\i,1.6) {\val};
    \node[anchor=east, font=\scriptsize, color=ForestGreen] at (-0.5,1.6) {prefix-min per block:};

    \foreach \i/\val in {0/2,1/2,2/2,3/2,4/3,5/3,6/3,7/6,8/1,9/4,10/4,11/7,12/2,13/3,14/3,15/3}
        \node[draw, minimum width=0.4cm, inner sep=0pt, font=\scriptsize\ttfamily, fill=orange!10] at (\i,0.4) {\val};
    \node[anchor=east, font=\scriptsize, color=OrangeRed] at (-0.5,0.4) {suffix-min per block:};

    \draw[rounded corners=1pt, color=red, ultra thick]
      (0.75, 2.5) rectangle (4.25, 3.5);
    \node[font=\tiny, color=red, anchor=south west] at (0.75, 3.55)
      {window 1: positions 1--4};

    \draw[->, color=red, thick] (1, 0.1) -- (1, -0.6);
    \draw[->, color=red, thick] (5, 1.3) -- (5, -0.6);
    \node[font=\scriptsize, color=red, anchor=north] at (3, -0.7)
      {$\min(\text{suffix}[1]=2,\;\text{prefix}[5]=5)\;=\;2$};
  \end{tikzpicture}
  \caption{The two-stacks idea (illustration adapted from
    Groot~Koerkamp \& Martayan,
    SIMD-Minimizers~\cite{groot_koerkamp_simd_minimizers_2025},
    Figure~2). The $k$-mer hash stream is logically split into blocks
    of $w$ hashes each. After every $w$-iteration block boundary the
    suffix-min array is rebuilt in $O(w)$ on top of the
    incrementally-maintained prefix-min array. Any window of size
    $w$ that overlaps two blocks splits cleanly into a suffix of one
    block and a prefix of the next -- so the window minimum is the
    minimum of \emph{exactly two} values from the two arrays. The
    one branch in the algorithm is the once-per-$w$-iters
    suffix-rebuild, which fires uniformly for all $L$ SIMD lanes.}
  \label{fig:two-stacks}
\end{figure}

The branch frequency drop is the headline win: V6 fires the rescan
on $\sim 34\,\%$ of iterations (data-dependent across four lanes);
V8 fires the suffix-rebuild on exactly every $w$-th iteration
($10\,\%$, data-independent).  Even the absolute branch count is
lower; the predictability is the bigger win because lanes no longer
diverge and the OoO engine can prefetch consistently.

\subsection{Per-step state in the V8 collector}

Algorithm~\ref{alg:two-stack} sketches the V8 inner loop. We omit
the $w$-rebuild body (it is a compile-time-unrolled loop of
\texttt{IfThenElse} pairs when $\texttt{StaticW}=10$) and the
strand-handling LSB-pack trick (it merges what would have been two
ring buffers into one by stuffing the strand bit into the lowest
bit of the hash; see~\cite[\S V4]{bioseq_eng_report_v2}).

\begin{algorithm}[H]
\caption{V8 \texttt{collect\_minimizers\_two\_stack\_simd} (per-step body)}
\label{alg:two-stack}
\small
\begin{algorithmic}[1]
\State \textbf{state}: $\textit{prefix\_min}[L]$, $\textit{ring}[w][L]$, $\textit{step}$, $\textit{ring\_idx}$
\State $\textit{hash} \gets \texttt{rolling\_kmer\_hash}.\textsc{Next}(d, c_\text{in}, c_\text{out})$
\State $\textit{strand}, \textit{palindrome} \gets \textsc{StrandMask}(\textit{hash})$
\State $\textit{valid} \gets \neg\textit{palindrome}$
\State $\textit{hs} \gets \textsc{IfThenElse}(\textit{strand}, \textit{hash}\,|\,1, \textit{hash}\,\&\,{\sim}1)$ \Comment{strand bit packed in LSB}
\State $\textit{elem} \gets \textsc{IfThenElse}(\textit{valid}, ((\textit{hash}\,\&\,\textbf{0xffff\_0000})\,|\,\textit{pos\_local}), \textbf{0xffff\_ffff})$
\State $\textit{ring}[\textit{ring\_idx}] \gets \textit{elem}$ \Comment{single ring buffer slot per lane}
\State $\textit{better} \gets \textit{elem} < \textit{prefix\_min}$
\State $\textit{prefix\_min} \gets \textsc{IfThenElse}(\textit{better}, \textit{elem}, \textit{prefix\_min})$
\State $\textit{ring\_idx} \mathrel{+}= 1$, $\textit{step} \mathrel{+}= 1$, $\textit{pos\_local} \mathrel{+}= 1$
\If{$\textit{ring\_idx} = w$}                                          \Comment{rebuild suffix}
   \State $\textit{ring}[i] \gets \min(\textit{ring}[i..w-1])$ \textbf{for all} $i$ \Comment{$O(w)$ uniform branch}
   \State $\textit{prefix\_min} \gets \textit{elem}$
   \State $\textit{ring\_idx} \gets 0$
\EndIf
\State $\textit{window\_min} \gets \min(\textit{prefix\_min}, \textit{ring}[\textit{ring\_idx}])$
\If{$\textit{step} \ge k+w-1$ \textbf{and} $\textit{window\_min}.\text{pos} \ne \textit{last\_emit\_pos}$}
   \State \textbf{emit} window minimum to per-lane scalar collector
   \State $\textit{last\_emit\_pos} \gets \textit{window\_min}.\text{pos}$
\EndIf
\end{algorithmic}
\end{algorithm}

The V8 numbers were already encouraging:
\texttt{mm\_sketch\_simd} dropped from V6's 24\,ms to 17\,ms at
903\,kbp -- a 30\,\% reduction on the minimizer-only phase, which
translated into the SIMD pipeline being only 1.4\,\% behind scalar
overall (down from V6's 8\,\% behind).  But ``almost as fast as
scalar'' is not ``10\,\% faster than scalar'', so we kept iterating.

\subsection{V8.5 -- defer the seven-stage scramble}
\label{sec:v8-defer-scramble}

The deepest lesson from V6's post-mortem -- ``the seven-stage hash
mix dominates the inner loop'' -- still applies to V8.  Every
position pays the seven dependent SIMD operations of
\texttt{hash64}, even though only $\sim 1/w$ of positions are
emitted as minimizers.

V8.5's idea is to defer the seven-stage scramble until emit time, and
only use a cheaper hash function for \emph{ordering} the values
inside the inner loop.  The cheap inner hash needs only one
property: be a bijection on the $2k$-bit kmer space, so that
$\textit{lex}(\textit{kmer}_1) = \textit{lex}(\textit{kmer}_2) \iff
\textit{kmer}_1 = \textit{kmer}_2$.  We use the
``lexicographic-canonical hash'' -- the bit-packed
$\min(\textit{forward}, \textit{rev\_compl})$ -- post-multiplied
by Knuth's golden-ratio constant
$0x9E3779B9$~\cite{knuth_taocp_v3}.  At emit time we run the
seven-stage \texttt{hash64} scramble on the chosen lane's
\emph{lexic} value (recovered from a parallel ring buffer), so the
emitted hash is bit-identical to what scalar minimap2 produces.
Both query and target apply the identical scheme, so the
minimizer-set is consistent across the two sides of the alignment.

The arithmetic of the saving:
\[
\underbrace{2.25 \cdot 10^5}_{\text{iters / Mbp}}
\;\times\;
\underbrace{(7\,\text{stages} - 2\,\text{stages})}_{\text{ops saved per iter}}
\;\approx\;
1.1 \cdot 10^6\,\text{SIMD ops saved}
\]
versus the
$\sim 22\,500$ \emph{additional} scalar scrambles paid at emit time
($\sim 30\,\text{cycles}$ each, total $\approx 0.21$\,ms).

The measured saving was $\sim 5$\,ms per 1\,Mbp pair on
\texttt{mm\_sketch\_simd}, taking the cost from V8's 17\,ms down to
12\,ms.  More importantly it cleared the way for the dual-stream and
inlining work that follow.

\section{Round 3 -- dual-stream interleaving (V11) and forced inlining (V12)}
\label{sec:v11-v12}

After V8.5, the V13-state phase breakdown
(Table~\ref{tab:v6-phase}'s rightmost column, computed from the
\texttt{report-long.json}) showed
\texttt{mm\_sketch\_simd} sitting at $\sim 13$\,ms -- still
slightly slower than scalar's 12.5\,ms.  We still had to find
$\sim 4$\,ms.  Two structural changes did it.

\subsection{V11 -- dual-stream interleaving}

A close reading of the assembly emitted by Apple LLVM revealed that
our 4-lane NEON inner loop has a 5-deep dependency chain per
iteration: \texttt{shift\_fwd $\to$ or $\to$ and $\to$ Min(fwd, rc)
$\to$ Mul(lex, golden)}.  The Apple\,M1 has \emph{four} 128-bit NEON
pipes that issue independently, so a single 4-lane stream wastes
three out of four pipes for any given cycle of the dependency
chain.

The fix is to keep \emph{two} independent 4-lane streams in flight
simultaneously and interleave their operations in source order:

\begin{verbatim}
shift_fwd_lo, shift_fwd_hi = ShiftLeft(state_fwd_lo / hi)
state_fwd_lo, state_fwd_hi = And(Or(shift_fwd_, bp_), kmer_mask)
state_rc_lo,  state_rc_hi  = Or(ShiftRight(state_rc_), new_high_)
lex_lo,       lex_hi       = Min(state_fwd_, state_rc_)
hash_lo,      hash_hi      = And(Mul(lex_, golden), hash_mask)
elem_lo,      elem_hi      = IfThenElse(valid_, hash_|pos_, max)
\end{verbatim}

The compiler does not need to discover the ILP; the source code
already exposes it. Each pair of operations operates on
register-disjoint values, so the M1 dispatch can issue both halves of
the same row in the same cycle.

We chose dual-stream rather than quad-stream because the per-stream
state plus the shared ring buffer for $w=10$ already consumes
$\sim 26$ NEON registers per stream; doubling to 52 with two streams
just barely fits inside the M1's 32-register file (the ring buffer
spills to stack, but the per-iteration state stays in registers --
which is exactly the right trade).  Going to four streams of 2-lane
vectors would spill the per-iteration state too, losing the ILP
benefit.

The wall-clock effect at 903\,kbp was a $\sim 4$\,ms reduction on
\texttt{mm\_sketch\_simd} (from 12\,ms to 8\,ms);
Figure~\ref{fig:phase-evolution} plots the full evolution.

\subsection{V12 -- the lambda-inline failure}

V11 was good but disappointing: the wall-clock saving was less than
the theoretical $2\times$ ILP win we expected.
\texttt{xcrun llvm-objdump -d build/bioseq/bioseq\_bench} on the V11
wrapper revealed the cause:

\begin{verbatim}
ldp q0, q1, [sp, #0x150]       # load bp_row from stack
ushr.4s v2, v1, #0x2           # extract per-BP value
movi.4s v1, #0x3
stp q0, q2, [sp, #0x150]       # store back
and.16b v0, v0, v1
and.16b v1, v2, v1
add x0, sp, #0x5, lsl #12      # prepare lambda's captured-this
add x0, x0, #0x770
movi.2d v2, #0xffffffffffffffff
add x1, x22, #0x8
bl 0x1002acca4 <step_one_bp_pair lambda>   # the offending call!
\end{verbatim}

Despite the per-step body being decorated with
\texttt{[[gnu::always\_inline]]}, Clang lowered the lambda as a
function call at every BP step.  Each call paid:

\begin{itemize}
\item Function-call prologue/epilogue (\texttt{stp x28-x30 [sp,...]}).
\item Captured-state spill/reload via stack -- every per-stream
  state register was loaded from the captured-this pointer at
  the start of each lambda invocation, and stored back at the
  end.
\item No ILP across iterations -- each lambda call is its own
  function frame, so dependent state spills to memory between
  consecutive iterations.
\end{itemize}

This is a known limitation of \texttt{[[gnu::always\_inline]]} on
captured-by-reference lambdas: the address-taken state escapes the
lambda body, so the compiler emits a real function call to be
defensive.  V12's fix is brutal but effective: rewrite the per-step
body as a \emph{free function template}
(\texttt{step\_one\_bp\_pair\_v3}) and pass all per-stream state by
reference.  Free functions reliably honour
\texttt{[[gnu::always\_inline]]}, and Clang does not treat
reference parameters as escaping when the function is fully
inlined.  Disassembly verification of V12 shows the per-step body
fully inlined into the function frame; the per-iter cost dropped
from $\sim 22$ SIMD ops + 1 \texttt{bl} call (the \texttt{bl}
having been $\sim 10\times$ more expensive than a single ALU op due
to the captured-state round-trip) to $\sim 15$ SIMD ops with no
function call.

V12 is the production minimizer at the time of writing.  Combined
with V8.5's deferred scramble and V11's dual-stream interleaving, the
inner \texttt{mm\_sketch\_simd} cost is now 9.2\,ms at 903\,kbp -- down
from V6's 24\,ms.  The combined \texttt{mm\_sketch} wrapper plus
\texttt{mm\_sketch\_simd} is 14.99\,ms versus scalar's 12.5\,ms; the
wrapper accounts for $\sim 2.5$\,ms of unavoidable per-call overhead
(the per-pair scalar emit collector, the bit-packing of the
collected output, the post-emit scramble) that scalar's
\texttt{mm\_sketch} subsumes into its single inner loop.  See
Figure~\ref{fig:phase-evolution} for the full evolution.

\begin{figure}[ht]
  \centering
  \begin{CacheMeCode}{python, custom include command={\input{\robExtAddCachePathAndName{\robExtFinalHash.pgf}}}, set placeholder eval={__LINEWIDTH__}{\lenToCmNoUnit[in]{\linewidth}}}
    import matplotlib.pyplot as plt
    import matplotlib
    import numpy as np
    matplotlib.use("pgf")
    matplotlib.rcParams.update({
      "font.family": "serif",
      "font.serif": [],
      "text.usetex": True,
    })
    fig, ax = plt.subplots()
    fig.set_size_inches(__LINEWIDTH__, 0.55*__LINEWIDTH__, forward=True)
    iterations = ['V6\n(rescan)', 'V8\n(two-stack)', 'V8.5\n(defer\nscramble)',
                  'V11\n(dual\nstream)', 'V12\n(inline)', 'scalar\n(reference)']
    mm_sketch_simd = [23.88, 19.90, 13.60, 13.60, 9.23, 0.0]
    mm_sketch_wrapper = [14.94, 12.43, 8.62, 8.62, 5.76, 12.50]
    sketch_total = [a+b for a,b in zip(mm_sketch_simd, mm_sketch_wrapper)]
    x = np.arange(len(iterations))
    ax.bar(x, mm_sketch_wrapper, label='mm\\_sketch (wrapper)', color='#9999cc')
    ax.bar(x, mm_sketch_simd, bottom=mm_sketch_wrapper,
           label='mm\\_sketch\\_simd (inner)', color='#cc4444')
    ax.axhline(12.50, linestyle='--', linewidth=1, color='#444444',
               label='scalar mm\\_sketch')
    for i, total in enumerate(sketch_total):
        ax.text(i, total + 0.5, f'{total:.1f}', ha='center', fontsize=7)
    ax.set_xticks(x)
    ax.set_xticklabels(iterations)
    ax.set_ylabel('Per-pair phase cost (ms) at 903\\,kbp', fontsize=8)
    ax.set_title('Minimizer-phase cost evolution across iterations', fontsize=9)
    ax.tick_params(axis='both', labelsize=7)
    ax.legend(fontsize=7, frameon=False, loc='upper right')
    ax.grid(True, alpha=0.3, axis='y')
    print(get_filename_from_extension(".pgf"))
    plt.savefig(get_filename_from_extension(".pgf"), bbox_inches="tight")
  \end{CacheMeCode}
  \caption{Evolution of the minimizer-phase cost
    (\texttt{mm\_sketch} wrapper plus \texttt{mm\_sketch\_simd}
    inner) at 903\,kbp across the iteration rounds discussed in
    \S\ref{sec:v6-rescan}--\S\ref{sec:v11-v12}. Numbers are means
    of three runs from the
    \texttt{[quick-long-sequence-benchmark-two-stack]} case at
    \texttt{-O3 -DNDEBUG}.  V8 is the two-stack port,
    V8.5 defers the minimap2 hash to emit time, V11 adds dual-stream
    interleaving, and V12 forces inlining of the per-step body via a
    free-function template.  The dashed horizontal line is the
    scalar \texttt{mm\_sketch} baseline.}
  \label{fig:phase-evolution}
\end{figure}

\section{Round 4 -- experiments that didn't work}
\label{sec:negative}

The optimisation work also produced several confidently-formulated
hypotheses that the data refused to confirm.  Documenting them is at
least as useful as documenting the wins, because it tells future
engineers what \emph{not} to spend time on.  In each case we built
the optimisation, ran the bench, and reverted.

\subsection{V14 -- true ntHash via NEON \texttt{TableLookupLanes}}
\label{sec:v14-nthash}

The V13 phase breakdown showed \texttt{collect\_seed\_hits} sitting
at 12\,ms for V12 versus 6.8\,ms for scalar -- a $\sim 5$\,ms
regression we attributed to the V8.5 deferred-scramble path
selecting AT-rich kmers more often than the
post-scramble-ordered scalar path.  AT-rich kmers naturally have
higher multiplicity in the genome, so they bloat the
\texttt{target\_index} hash buckets and the per-query lookup chain.

The clean fix would have been to use ntHash~\cite{nthash_2016}, the
rolling hash recommended by the SIMD-Minimizers
paper~\cite{groot_koerkamp_simd_minimizers_2025}.  ntHash is a
sequence of \texttt{rotate\,+\,XOR\,+\,table-lookup} operations
where the table is indexed by the 2-bit nucleotide rank, so on NEON
it can be implemented as a $4$-entry \texttt{tbl.16b} lookup per
character.  This produces a uniform hash distribution across the
$k$-mer space, which removes the AT-bias.

V14 implements canonical ntHash inside the V13 dual-stream chassis.
The implementation is correct -- it preserves the warmup-cancellation
invariant by precomputing constants that subtract the
$c_\text{old} = 0$-padding contribution after $k$ rolling steps --
and the integration suite passes.

What V14 \emph{didn't} buy was wall-clock:

\begin{itemize}
\item \texttt{mm\_sketch\_simd} cost \emph{increased} by 2\,ms
  (the table lookup adds a NEON \texttt{tbl.16b}, two rotates, and
  some XORs per stream per BP).
\item \texttt{collect\_seed\_hits} dropped by \emph{0.4\,ms}, not
  the predicted 5\,ms.
\item Net: V14 is statistically tied with V13 at LONG and
  VERY\_LONG, and slightly slower at XL.
\end{itemize}

The lesson is that the AT-bias hypothesis was wrong.  The
\texttt{collect\_seed\_hits} regression is dominated not by the
minimizer hash distribution but by the \texttt{std::unordered\_map}
lookup overhead: bioseq's target index uses STL's default hash
table, while the C minimap2 path uses an open-addressing
\texttt{khash\_t<uint64\_t>}.  The next-most-impactful single edit
for closing Goal~2 -- replacing \texttt{std::unordered\_map} with a
flat hash on a uint32 $\to$ vector layout -- is in the post-V14
future-work list of the engineering report and is not part of this
chapter.

\subsection{V15 -- dual-batch chaining}
\label{sec:v15-chain-v2}

The dual-stream interleaving idea that delivered for the minimizer
loop (\S\ref{sec:v11-v12}) also worked beautifully for an aggressive
unrolling of the SIMD chaining inner loop. Unfortunately, the
chaining loop is not register-bound and 4-pipe-issue-limited;
\emph{it is load-store-bound}.  V15 doubles the SoA loads per inner
iteration from 4 to 8, which queues up on the M1's 4 load-store
pipes and stalls the issue queue.  Net wall-clock effect:

\begin{itemize}
\item \texttt{mg\_lchain\_simd} regresses by 5--13\,\% at 903\,kbp.
\item Total pipeline regresses by 7\,\% versus V13 chaining + V13 minimizer.
\end{itemize}

The lesson is that the unrolling-with-multiple-reduction-streams
pattern only helps register-bound, issue-width-limited kernels.
Apply it only after confirming the target kernel is issue-width
limited, not load-store limited.  V15 stays in the codebase as a
documented negative result.

\subsection{Goal~2-style \texttt{prefix\_kmer} drop}

A future-work item from V12 proposed dropping the parallel
\texttt{ring\_buf\_hs} buffer and the \texttt{prefix\_kmer} state
register by packing the strand bit into bit~16 of \texttt{elem}.
The estimated saving was $\sim 3$ ops/stream/iter and 22 vec\_t of
register-pressure relief.  We implemented and benchmarked.  Result:

\begin{verbatim}
                                  V12 (pre)   V12 (post-Goal-2)   delta
mm_sketch_simd                    9.0 ms      9.5 ms              +0.5 ms (matches estimate)
target_index_build                24.7 ms     4.9 ms              -19.8 ms (~85% fewer keys)
collect_seed_hits                 12.3 ms     842.0 ms            +830 ms (70x regression)
mg_lchain_simd                    8.9 ms      299.9 ms            +291 ms (34x regression)
mm_align_skeleton                 87.2 ms    4811.0 ms            +4724 ms (56x regression)
\end{verbatim}

The kmer space stored in \texttt{elem} after the trick was only 14
bits -- not the 16 bits the V12 future-work entry assumed -- so the
\texttt{target\_index} hash buckets collapsed to $\sim 16\,000$
distinct keys for $168\,000$ query minimizers.  Each query had
$\sim 10$ collisions per bucket, the chain-anchor count exploded,
and the alignment skeleton had to walk through all of them.  The
pipeline is $36\times$ slower as a result.

The lesson is that hash-precision losses in a sketching scheme
\emph{compound} catastrophically through the downstream pipeline.
Even an apparently small precision loss must be validated against
the full pipeline's wall-clock, not just the sketch phase's.

\section{Round 5 -- the divergence-estimator fix (V13)}
\label{sec:v13-est-err}

The single biggest \emph{non-minimizer} contribution to closing
Goal~2 is the V13 rewrite of
\texttt{estimate\_divergence\_from\_seeds} in
\path{bioseq/backend/sequences/minimap_port.cppm}, which dropped
\texttt{mm\_est\_err} from V12-pre's 7.88\,ms to V13's 0.85\,ms --
a $9\times$ speedup, removing one of the two large downstream
regressions in the SIMD pipeline.

\subsection{Why \texttt{mm\_est\_err} mattered}

For each chain, minimap2 estimates the per-base sequence divergence
from a Poisson model on the minimizer-set
intersection~\cite{minimap2_paper}:
\[
\hat{\epsilon} = \frac{1}{k} \log \frac{n}{m}
\]
where $n$ is the number of query minimizers in the chain's reference
span and $m$ is the number of those that are actually seeded
(i.e.\ that appear as anchors in the chain).  Computing $n$ requires
\emph{knowing which positions are query minimizers}, which is a
membership test against the query's minimizer-position set.

The pre-V13 implementation:

\begin{enumerate}
\item Built a sorted array \texttt{minimizer\_positions[]} of every
  query minimizer's position via \texttt{std::sort}.  The V12
  dual-stream collector emits in lane-interleaved order, so the
  array is \emph{not} sorted on output, and \texttt{std::sort}
  sees worst-case input.
\item For each chain anchor, did a \texttt{std::lower\_bound} into
  the sorted array to test membership.
\end{enumerate}

The dominant cost was the \texttt{std::sort} on 168\,k entries:
$\sim 5$--$6$\,ms per pair.  The \texttt{lower\_bound} loop added
another $\sim 1$--$2$\,ms.

\subsection{The sort-free prefix-count rewrite}

V13 replaces both with a single linear-pass prefix-sum:

\[
\textit{prefix\_count}[i] = \big|\{ \, p \in \textit{minimizer\_positions} \mid p < i \,\}\big|
\]

Built in $O(N + L)$ where $N$ is the minimizer count and $L$ is the
query length.  After construction:

\begin{itemize}
\item $\textit{is\_minimizer}(p) = (\textit{prefix\_count}[p+1] > \textit{prefix\_count}[p])$ -- $O(1)$.
\item $\textit{range\_count}(a, b) = \textit{prefix\_count}[b+1] - \textit{prefix\_count}[a]$ -- $O(1)$.
\end{itemize}

The per-region inner loop is now a single $O(\textit{seed\_count})$
walk via $O(1)$ probes.  No sort, no binary search.

The fix landed for \emph{all} configs (scalar pipeline included) so
even the scalar path got a small speedup: scalar
\texttt{mm\_est\_err} dropped from 1.21\,ms to 0.85\,ms, V12-SIMD
\texttt{mm\_est\_err} dropped from 7.88\,ms to 0.85\,ms.  The
SIMD-minus-scalar regression is gone.

There was also a correctness bug exposed by the rewrite: the V10
``reverse-chain fix'' was wrong, because for both forward and
reverse chains the lambda \texttt{compute\_forward\_pos} was
already producing monotonically-increasing forward positions in
chain order. We removed the reverse-chain split entirely.

\subsection{Net pipeline impact}

Figure~\ref{fig:phase-stack-v13} stacks the per-phase cost at
903\,kbp for V12-pre versus V13.  The two large gains are
\texttt{mm\_est\_err} (from 7.88\,ms to 0.85\,ms) and the V12
\texttt{mm\_sketch} stack (from 14.99\,ms to 9.23\,ms, already
covered in \S\ref{sec:v11-v12}).

\begin{figure}[ht]
  \centering
  \begin{CacheMeCode}{python, custom include command={\input{\robExtAddCachePathAndName{\robExtFinalHash.pgf}}}, set placeholder eval={__LINEWIDTH__}{\lenToCmNoUnit[in]{\linewidth}}}
    import matplotlib.pyplot as plt
    import matplotlib
    import numpy as np
    matplotlib.use("pgf")
    matplotlib.rcParams.update({
      "font.family": "serif",
      "font.serif": [],
      "text.usetex": True,
    })
    fig, ax = plt.subplots()
    fig.set_size_inches(__LINEWIDTH__, 0.55*__LINEWIDTH__, forward=True)
    configs = ['scalar', 'V6 SIMD', 'V8 two-stack', 'V12 v3 (pre-V13)', 'V12 v3 (post-V13)']
    phases = [
        ('mm\\_sketch + mm\\_sketch\\_simd', [12.50, 38.82, 32.43, 14.99, 14.52], '#9999cc'),
        ('target\\_index\\_build',         [16.0,  14.3,  14.6,  26.38, 23.83], '#1f77b4'),
        ('collect\\_seed\\_hits',          [6.5,   11.0,  11.3,  12.32, 12.0],  '#2ca02c'),
        ('mg\\_lchain\\_dp/simd',          [16.0,  8.9,   8.9,   8.88,  8.78],  '#ff7f0e'),
        ('mm\\_est\\_err',                 [0.7,   7.2,   7.1,   7.88,  0.85],  '#d62728'),
        ('mm\\_align\\_skeleton',          [87.5,  86.4,  86.5,  87.25, 85.42], '#888888'),
    ]
    bottoms = np.zeros(len(configs))
    for name, vals, color in phases:
        ax.bar(configs, vals, bottom=bottoms, label=name, color=color)
        bottoms = bottoms + np.array(vals)
    for i, b in enumerate(bottoms):
        ax.text(i, b + 1.5, f'{b:.0f}', ha='center', fontsize=7)
    ax.set_ylabel('Per-pair phase cost (ms) at 903\\,kbp', fontsize=8)
    ax.set_title('Per-phase pipeline composition (903\\,kbp, mean of 3 runs)', fontsize=9)
    ax.tick_params(axis='both', labelsize=7)
    ax.legend(fontsize=6, frameon=False, loc='upper right', ncol=2)
    ax.grid(True, alpha=0.3, axis='y')
    plt.xticks(rotation=15, ha='right', fontsize=7)
    print(get_filename_from_extension(".pgf"))
    plt.savefig(get_filename_from_extension(".pgf"), bbox_inches="tight")
  \end{CacheMeCode}
  \caption{Per-phase composition of the per-pair pipeline at
    903\,kbp for the headline configurations.  ``V12 v3 (pre-V13)''
    is the V12 collector with the pre-V13 \texttt{mm\_est\_err}
    code; ``V12 v3 (post-V13)'' is the same minimizer collector with
    the V13 sort-free prefix-count rewrite.  The pre/post pair lets
    the reader visually separate the V12 (minimizer-side) wins from
    the V13 (\texttt{mm\_est\_err}) wins.  The \texttt{mm\_align\_skeleton}
    bar is identical across all configs, accounting for $\sim 50\,\%$
    of total wall-clock, and is not addressed by any of the work in
    this chapter.}
  \label{fig:phase-stack-v13}
\end{figure}

\section{Final results}
\label{sec:final-results}

The combination of the V12 minimizer collector and the V13
\texttt{mm\_est\_err} rewrite is the production state at the time of
writing.  Across the five sequence-length buckets the bench
exercises (LONG, VERY\_LONG, XL), V13 is the fastest of the SIMD
configurations on every bucket and is the fastest config
\emph{overall} on the LONG sub-buckets (47\,kbp, 49\,kbp, 52\,kbp).
At VERY\_LONG (903\,kbp) it is 8.8\,\% faster than scalar -- just
short of the 10\,\% Goal-2 target.  At XL (5\,Mbp) it is 4.9\,\%
faster than scalar.

\begin{figure}[ht]
  \centering
  \begin{CacheMeCode}{python, custom include command={\input{\robExtAddCachePathAndName{\robExtFinalHash.pgf}}}, set placeholder eval={__LINEWIDTH__}{\lenToCmNoUnit[in]{\linewidth}}}
    import matplotlib.pyplot as plt
    import matplotlib
    import numpy as np
    matplotlib.use("pgf")
    matplotlib.rcParams.update({
      "font.family": "serif",
      "font.serif": [],
      "text.usetex": True,
    })
    fig, ax = plt.subplots()
    fig.set_size_inches(__LINEWIDTH__, 0.50*__LINEWIDTH__, forward=True)
    buckets = ['47 kbp', '49 kbp', '52 kbp', '903 kbp', '5.0 Mbp']
    delta_v6   = [(15.21 - 14.56) / 14.56 * 100,
                  (15.90 - 14.94) / 14.94 * 100,
                  (16.05 - 15.09) / 15.09 * 100,
                  (201.68 - 188.93) / 188.93 * 100,
                  (1207.99 - 1171.0) / 1171.0 * 100]
    delta_v8   = [(14.64 - 14.56) / 14.56 * 100,
                  (15.07 - 14.94) / 14.94 * 100,
                  (15.24 - 15.09) / 15.09 * 100,
                  (183.77 - 188.93) / 188.93 * 100,
                  (1131.87 - 1171.0) / 1171.0 * 100]
    delta_v13  = [(13.78 - 14.56) / 14.56 * 100,
                  (14.02 - 14.94) / 14.94 * 100,
                  (13.82 - 15.09) / 15.09 * 100,
                  (172.38 - 188.93) / 188.93 * 100,
                  (1113.3 - 1171.0) / 1171.0 * 100]
    x = np.arange(len(buckets))
    width = 0.27
    ax.bar(x - width, delta_v6,  width, label='V6 SIMD (rescan)', color='#1f77b4')
    ax.bar(x,         delta_v8,  width, label='V8 two-stack',     color='#2ca02c')
    ax.bar(x + width, delta_v13, width, label='V13 (production)', color='#d62728')
    ax.axhline(-10, linestyle='--', linewidth=1, color='#ff0000',
               label='-10\\% Goal-2 target')
    ax.axhline(0, linewidth=0.6, color='#000000')
    ax.set_xticks(x)
    ax.set_xticklabels(buckets)
    ax.set_ylabel('Wall-clock delta vs scalar (\\%, negative = faster)', fontsize=8)
    ax.set_title('SIMD vs scalar speedup across iterations', fontsize=9)
    ax.tick_params(axis='both', labelsize=7)
    ax.legend(fontsize=7, frameon=False, loc='upper right', ncol=2)
    ax.grid(True, alpha=0.3, axis='y')
    print(get_filename_from_extension(".pgf"))
    plt.savefig(get_filename_from_extension(".pgf"), bbox_inches="tight")
  \end{CacheMeCode}
  \caption{Wall-clock delta versus scalar for the three headline
    SIMD configurations, expressed as percent (negative = faster).
    The dashed line is the $-10\,\%$ Goal-2 target. V13 approaches
    the target across the LONG / VERY\_LONG buckets (between
    $-8.4\,\%$ and $-8.8\,\%$ on the plotted three-run means) and
    is within 5\,\% of the target on XL.}
  \label{fig:speedup-summary}
\end{figure}

The summary table:

\begin{table}[ht]
  \centering
  \small
  \begin{tabular}{lrrrr}
    \toprule
    Length     & scalar (ms) & V13 (ms) & V13 vs scalar & V13 vs V6 SIMD \\
    \midrule
    47\,kbp   &  14.56  & 13.78 & $-5.4\,\%$ & $-9.4\,\%$ \\
    49\,kbp   &  14.94  & 14.02 & $-6.2\,\%$ & $-11.8\,\%$ \\
    52\,kbp   &  15.09  & 13.82 & $-8.4\,\%$ & $-13.9\,\%$ \\
    903\,kbp  & 188.93  & 172.38 & $-8.8\,\%$ & $-14.5\,\%$ \\
    5.0\,Mbp  & 1171.0  & 1113.3 & $-4.9\,\%$ & $-7.8\,\%$ \\
    \bottomrule
  \end{tabular}
  \caption{Final V13-state speedup numbers across the five
    sequence-length buckets, expressed as wall-clock per pair
    (mean of 3 runs from
    \texttt{[quick-long-sequence-benchmark]}, bench-opt
    \texttt{-O3 -DNDEBUG}, CCS error rate, seed
    \texttt{0xBEEFCAFE12345678}). The $-10\,\%$ Goal-2 target is
    narrowly missed at 903\,kbp and 5.0\,Mbp; on the LONG sub-buckets
    (47, 49, 52\,kbp) it is met or comfortably approached.}
  \label{tab:final-numbers}
\end{table}

\section{What we learned about SIMD bioinformatics}
\label{sec:lessons}

Several lessons generalise beyond the minimizer collector:

\begin{enumerate}
\item \textbf{Data-dependent branches are exponentially more
  expensive in SIMD than in scalar.}  In scalar code, the OoO
  engine's branch predictor has $\sim 90\,\%$ accuracy on a
  $1/w$-rate branch.  In a 4-lane SIMD code, ``no lane mispredicts''
  is $0.9^4 \approx 65\,\%$.  The data-independent two-stacks
  algorithm replaces the worst-case branching inner loop with a
  uniform every-$w$-iters branch, which costs nothing.

\item \textbf{Hash-function deferral is a free win for any sketching
  scheme.}  Whenever a cheap order-preserving hash exists, the
  expensive ``real'' hash should be applied at emit time on the
  $\sim 1/w$ subset of positions actually emitted, not in the
  inner loop.

\item \textbf{Lambdas are not always inlined, even with
  \texttt{[[gnu::always\_inline]]}.}  Captured-by-reference lambdas
  with non-trivial bodies are at risk.  Free-function templates
  with state passed by reference are reliable.  The disassembly is
  the only way to confirm.

\item \textbf{Two register-bound 4-lane streams beat one
  register-bound 4-lane stream on M1.}  The chassis is a
  register-bound, issue-width-limited inner loop; doubling the
  parallelism halves the steady-state critical-path latency.  Going
  to four streams overflows the register file.

\item \textbf{Don't apply the dual-stream pattern to
  load-store-bound kernels.}  V15's chaining V2 result is
  unambiguous: the M1 has only 4 load-store pipes, and a kernel that
  is already issuing more than 4 loads per cycle does not benefit
  from doubling them.

\item \textbf{Hash-precision losses compound through the pipeline.}
  Even a 2-bit precision loss in the inner ordering hash can produce
  a $36\times$ end-to-end slowdown (the Goal-2 \texttt{prefix\_kmer}
  drop result, \S\ref{sec:negative}). Validate every change against
  the full pipeline, not just the sketch phase.

\item \textbf{The minimizer phase is not the only phase that pays
  for SIMD.}  When the SIMD chainer produces a different anchor
  distribution than scalar, downstream stages
  (\texttt{collect\_seed\_hits}, \texttt{mm\_est\_err}) regress.
  Some of those regressions are structural (different chain shape)
  and require restructuring those stages -- e.g.\ the V13
  sort-free \texttt{mm\_est\_err} rewrite -- not just smarter SIMD
  in the sketch.

\item \textbf{The alignment skeleton is the unmovable mass.}
  At 903\,kbp, \texttt{mm\_align\_skeleton} accounts for $\sim 50\,\%$
  of the pipeline wall-clock, and it is identical across all
  configurations.  Cutting the entire minimizer + chaining +
  est\_err + seed\_hits stack to zero would buy at most $50\,\%$
  speedup.  Goals like ``-10\,\%'' on the SIMD pipeline are very
  achievable; Goals like ``2$\times$'' would require also moving
  the alignment phase to SIMD (a substantial KSW2 rewrite or a
  port to Block-Aligner~\cite{block_aligner_paper}).
\end{enumerate}

\section{Future work}
\label{sec:future-work}

The post-V13 future-work list, in approximate priority order:

\begin{enumerate}
\item \textbf{Replace \texttt{std::unordered\_map} for the target
  index.}  An open-addressing flat hash map (e.g.\
  \texttt{absl::flat\_hash\_map} or a hand-rolled
  \texttt{khash\_t}-style probing table) is the V14-confirmed
  largest-single-edit win for closing
  \texttt{collect\_seed\_hits}'s 5\,ms regression.  This is
  not minimizer code at all -- it is the index data structure --
  but it is the fastest path to clearing Goal-2 on the LONG and
  VERY\_LONG buckets robustly.
\item \textbf{Switch the bench storage to \texttt{uint64\_t}.}
  Halves the SIMD lane count (4 $\to$ 2 on NEON) but eliminates the
  \texttt{prefix\_kmer} ring buffer entirely while preserving full
  kmer precision.  Whether the lane-count loss outweighs the
  register-pressure relief is an empirical question.
\item \textbf{Sorted minimizer emission in the V12 collector.}  If
  the per-lane outputs are merged at the end of the SIMD body
  rather than in lane-interleaved order, downstream stages that
  rely on sorted input (notably \texttt{target\_index\_build}) get
  a free $\sim 1$\,ms win.
\item \textbf{Lazy \texttt{mm\_est\_err} bitmap.}  For pairs whose
  chain spans the full sequence, the per-region walk is
  $O(\text{seed\_count})$ anyway, and we could skip the
  prefix-count construction with an $O(1)$ heuristic.
\item \textbf{Batched KSW2 alignment.}  The unmovable mass.
  Substantial multi-day rewrite, but the only way to get above
  $\sim 12\,\%$ end-to-end speedup.
\end{enumerate}

The codebase is at the V13 state, with V14 and V15 retained as
documented negative-result research artefacts.  Researchers
extending this work should consult the engineering reports
(\texttt{ENGINEERING\_REPORT.md} for V3 and earlier,
\texttt{ENGINEERING\_REPORT\_V2.md} for V4 onwards) for the
per-iteration step-by-step record, and the \path{view\_report.ipynb}
notebook (Tier-X minimizer-variant comparison section) for the
auto-generated charts that mirror Figures~\ref{fig:headline-result},
\ref{fig:phase-evolution} and~\ref{fig:speedup-summary} above.
